\section{The inverse matrix method}

Some systems of linear equations can be solved by using the inverse of the coefficient matrix. This method is based on the matrix equation
\[
Ax=b.
\]
Here \(A\) is the coefficient matrix, \(x\) is the column vector of unknowns, and \(b\) is the right-hand side vector.

The method is simple in form, but its meaning should be understood carefully. If \(A\) is invertible, then multiplying by \(A^{-1}\) undoes the action of \(A\).

\subsection*{Solving \(Ax=b\)}

Suppose
\[
Ax=b
\]
and \(A\) is invertible. Multiply both sides on the left by \(A^{-1}\):
\[
A^{-1}Ax=A^{-1}b.
\]
Since
\[
A^{-1}A=I,
\]
we get
\[
Ix=A^{-1}b.
\]
Therefore
\[
x=A^{-1}b.
\]

\begin{theorembox}
If \(A\) is an invertible square matrix, then the system
\[
Ax=b
\]
has the unique solution
\[
x=A^{-1}b.
\]
\end{theorembox}

The word ``invertible'' is essential. If \(A^{-1}\) does not exist, this method cannot be used.

\subsection*{A worked example}

\begin{examplebox}
Solve
\[
\begin{cases}
2x+y=5,\\
x+y=3.
\end{cases}
\]
Write the system as
\[
\begin{pmatrix}
2&1\\
1&1
\end{pmatrix}
\begin{pmatrix}
x\\y
\end{pmatrix}
=
\begin{pmatrix}
5\\3
\end{pmatrix}.
\]
Thus
\[
A=\begin{pmatrix}2&1\\1&1\end{pmatrix},
\qquad
b=\begin{pmatrix}5\\3\end{pmatrix}.
\]
The determinant is
\[
\det A=2\cdot1-1\cdot1=1,
\]
so \(A\) is invertible. For a \(2\times2\) matrix,
\[
A^{-1}=\begin{pmatrix}1&-1\\-1&2\end{pmatrix}.
\]
Therefore
\[
\begin{pmatrix}x\\y\end{pmatrix}
=A^{-1}b
=
\begin{pmatrix}1&-1\\-1&2\end{pmatrix}
\begin{pmatrix}5\\3\end{pmatrix}.
\]
Compute the product:
\[
\begin{pmatrix}x\\y\end{pmatrix}
=
\begin{pmatrix}5-3\\-5+6\end{pmatrix}
=
\begin{pmatrix}2\\1\end{pmatrix}.
\]
Thus
\[
x=2,
\qquad
 y=1.
\]
\end{examplebox}

The inverse matrix method gives the same result as other correct methods. The difference is the language: here solving means applying the inverse transformation.

\subsection*{Why the method works}

The equation
\[
Ax=b
\]
can be read as follows: the matrix \(A\) acts on the unknown vector \(x\) and produces the vector \(b\). If \(A\) is invertible, then this action can be reversed. Applying \(A^{-1}\) to \(b\) recovers the original vector \(x\).

This interpretation is important. The formula
\[
x=A^{-1}b
\]
is not magic. It is the matrix version of undoing an operation.

\subsection*{Connection with determinants}

For a square matrix \(A\), invertibility is equivalent to
\[
\det A\ne0.
\]
Therefore the inverse matrix method applies exactly when the coefficient matrix has non-zero determinant.

If
\[
\det A=0,
\]
then \(A^{-1}\) does not exist. The system may have no solution or infinitely many solutions, but it cannot be solved by multiplying by \(A^{-1}\).

\begin{examplebox}
Consider
\[
\begin{cases}
x+y=2,\\
2x+2y=4.
\end{cases}
\]
The coefficient matrix is
\[
A=\begin{pmatrix}1&1\\2&2\end{pmatrix},
\]
and
\[
\det A=0.
\]
The inverse matrix method cannot be used. However, the system does have solutions: every pair \((x,y)\) satisfying \(x+y=2\) is a solution. There are infinitely many such pairs.
\end{examplebox}

Again, the failure of invertibility means loss of uniqueness, not automatically inconsistency.

\subsection*{When this method is useful}

The inverse matrix method is conceptually clear, but it is not always the most practical computational method. Finding \(A^{-1}\) can be more work than solving the system directly by elimination.

The method is especially useful when:
\begin{itemize}
\item the inverse matrix is already known,
\item several systems have the same coefficient matrix \(A\) but different right-hand sides,
\item we want to understand the system as a matrix equation.
\end{itemize}

For routine computation, Gaussian elimination is usually more flexible.

\begin{summarybox}
When using the inverse matrix method, ask:
\begin{itemize}
\item Can the system be written as \(Ax=b\)?
\item Is \(A\) square?
\item Is \(A\) invertible, equivalently \(\det A\ne0\)?
\item Have I multiplied on the correct side by \(A^{-1}\)?
\item Does the vector \(A^{-1}b\) satisfy the original system?
\end{itemize}
\end{summarybox}

The inverse matrix method is powerful because it connects solving systems with the idea of reversing a linear transformation.